\section{国内外研究现状}\label{sec:related-work}

机器人基础软件层面的研究可大致归为三个方向：操作系统、通信构件与开发框架；
本节按此三个方向梳理国内外现有工作，并结合本文关注的端到端推理主循环指出研究空白。

在操作系统层面，研究核心是实时性保证与资源调度。
标准 Linux 内核采用公平调度机制，能够兼顾通用计算负载，但难以直接满足机器人控制中的确定性时序需求。
PREEMPT-RT 补丁通过增强内核可抢占性改善系统响应能力，是机器人控制系统中常见的 Linux 实时化方案\cite{rtlinux}；
相关研究同时指出，在并发通信和多任务负载下，仅依靠实时补丁仍不足以稳定约束延迟，还需要结合 CPU 隔离、绑核和优先级配置等手段。
ROS2 实时性研究综述\cite{ros2rtsurvey} 进一步表明，执行器模型、回调调度器和内核调度策略是当前机器人实时系统研究的重要分支。
近期 AMS\cite{ams} 将面向动作执行的操作系统原语引入具身智能推理管理，
通过显式管理动作上下文、异常和重放过程提升了真实机械臂任务的执行可靠性。
这些工作说明，操作系统并不是机器人学习系统的透明底座，而会直接影响推理任务的稳定执行。
不过，现有研究多从实时调度机制或动作管理接口出发，较少进一步分析深度学习推理框架、Python 运行时和设备 I/O 在同一控制闭环中的耦合关系。

在通信构件层面，ROS2\cite{ros2} 以 DDS（Data Distribution Service）为底层通信机制，
已经成为学术界与工业界广泛采用的机器人中间件。
国内机器人系统研究中，ROS 在焊接机器人场景中的仿真平台搭建与任务规划能力得到了验证\cite{houguowei2023weldingrobot}，
在喷浆机器人场景中的系统设计与轨迹规划能力也得到了验证\cite{linxuebin2020shotcreterobot}。
这类工作强调机器人任务建模、轨迹规划和系统集成，使复杂机械臂应用能够在统一软件框架下完成验证。
与此同时，面向 ARM 嵌入式平台的研究发现，节点组合方式、进程边界和消息传递路径会显著影响 CPU 占用与通信延迟\cite{ros2composition}；
在树莓派等低功耗平台上，Linux 网络栈也可能成为硬实时通信的关键限制因素\cite{ros2ddsrpi}。
为支撑通信链路分析，ROS2 追踪工具提供了基于 LTTng 的低开销追踪方法\cite{ros2tracing}，
能够观察 ROS2 节点调度与消息传递过程。
不过，上述工作主要围绕机器人中间件和通信链路展开，
对于不完全依赖 ROS2、而是由学习框架直接驱动摄像头、舵机和神经网络推理的在线控制程序，其分析范围仍然不够完整。

在开发框架层面，深度学习推理框架已成为机器人基础软件的重要组成部分。
TFLite\cite{tflite} 面向资源极其受限的端侧场景，强调量化推理和轻量化部署；
ONNX Runtime\cite{onnxruntime} 通过面向 ARM64 的算子优化提高端侧 CPU 推理效率；
PyTorch\cite{pytorch} 则凭借灵活的模型表达和较完整的生态支持，成为 LeRobot 等机器人开源框架常用的推理后端。
在机器人学习算法方面，ACT、Diffusion Policy 等策略推动了模仿学习和端到端控制的发展，
HuggingFace LeRobot\cite{lerobot} 等开源框架也降低了策略训练、数据采集和部署复现的门槛。
然而，这类研究通常更关注模型精度、任务成功率和算法泛化能力，
对低成本 ARM 平台上在线推理时的系统级开销关注不足。
尤其是推理框架内部线程池、Python 运行时调度、Linux 内核调度和硬件访问路径之间并不存在统一的时序管理机制，
在资源受限平台上可能产生阻塞、等待和上下文切换等隐性开销，而这些开销很难仅通过模型层或算子层优化解释。

综上，已有工作分别从操作系统实时化、通信中间件和推理框架优化等角度推动了机器人基础软件的发展，
但研究重点大多停留在单一层面或单一链路。
操作系统研究强调调度与实时性，通信构件研究强调消息传递，中间件追踪工具主要覆盖 ROS2 通信路径，
推理框架研究则更重视模型执行效率和算子优化。
对于 LeRobot 这类由 Python 在线控制程序组织感知、推理和执行的机器人学习系统，
从 AI 推理框架到 Python 运行时、Linux 内核再到硬件 I/O 的端到端调用链仍缺乏系统性剖析。
尤其是系统调用阻塞、设备 I/O 等低层路径通常需要结合 strace、ftrace 与内核插桩等工具共同观测，
现有工作较少将这些追踪结果与机器人主循环的阶段耗时直接对应起来。
系统性能分析方法论强调，应从应用、运行时、操作系统和硬件路径联合定位瓶颈\cite{perfbook}。
因此，本文建立覆盖 AI 推理框架层、Python 运行时层、OS 内核层与硬件 I/O 层的四层性能模型，
并结合阶段计时、系统调用分析和内核追踪方法，
定位 ARM 嵌入式平台上机器人在线控制主循环的真实瓶颈。
